% APPENDIX: STATISTICAL ANALYSIS DETAILS
\label{sec:appendix-stats}

\subsection{Baseline Statistics}

Table~\ref{tab:baseline-stats-full} reports comprehensive summary statistics from the academic baseline experiment (N=100 runs, 2{,}000 steps each).

\begin{table}[htbp]
\centering
\caption{Baseline simulation summary statistics (N=100 runs)}
\label{tab:baseline-stats-full}
\small
\begin{tabular}{@{}lrrrr@{}}
\toprule
Metric & Mean & SD & 95\% CI & Range \\
\midrule
Proposal Pass Rate & 0.879 & 0.026 & [0.874, 0.884] & [0.81, 0.94] \\
Average Turnout & 0.348 & 0.021 & [0.344, 0.352] & [0.30, 0.40] \\
Quorum Reach Rate & 0.999 & 0.002 & [0.999, 1.000] & [0.99, 1.00] \\
Voter Participation & 0.304 & 0.012 & [0.301, 0.306] & [0.27, 0.33] \\
Token Gini & 0.409 & 0.043 & [0.400, 0.417] & [0.31, 0.52] \\
Whale Influence & 0.238 & 0.014 & [0.235, 0.241] & [0.20, 0.27] \\
Capture Risk & 0.226 & 0.013 & [0.223, 0.229] & [0.19, 0.26] \\
Delegate Conc. & 0.219 & 0.013 & [0.217, 0.222] & [0.19, 0.25] \\
Treasury Final & 11{,}446 & 583 & [11{,}331, 11{,}562] & [9{,}980, 12{,}890] \\
Total Proposals & 203.2 & 15.7 & [200.1, 206.3] & [165, 242] \\
\bottomrule
\end{tabular}
\end{table}

\subsubsection{Confidence Interval Computation}

Confidence intervals are computed using Student's $t$-distribution to account for estimation uncertainty in the sample standard deviation:
\begin{equation}
\text{CI}_{1-\alpha} = \bar{x} \pm t_{n-1, \alpha/2} \cdot \frac{s}{\sqrt{n}}
\end{equation}
where $\bar{x}$ is the sample mean, $s$ is the sample standard deviation, $n$ is the number of simulation runs, and $t_{n-1, \alpha/2}$ is the critical value from the $t$-distribution with $n-1$ degrees of freedom.

For our sample sizes ($n = 100$), $t_{99, 0.025} \approx 1.984$ for 95\% CIs.

\subsubsection{Power Analysis}

Table~\ref{tab:power-analysis} reports the statistical power achieved for each experiment based on sample size and assumed medium effect size ($d = 0.5$).

\begin{table}[htbp]
\centering
\caption{Statistical power analysis by experiment ($\alpha = 0.05$, target $d = 0.5$)}
\label{tab:power-analysis}
\small
\begin{tabular}{@{}lrrrl@{}}
\toprule
Experiment & N (per config) & Power & Min Detectable $d$ & Adequacy \\
\midrule
00 Baseline & 100 & $>$0.99 & 0.28 & Excellent \\
01 Calibration & 100 & $>$0.99 & 0.28 & Excellent \\
02 Ablation & 100 & $>$0.99 & 0.28 & Excellent \\
03 Quorum & 100 & $>$0.99 & 0.28 & Excellent \\
04 Capture & 100 & $>$0.99 & 0.28 & Excellent \\
05 Pipeline & 100 & $>$0.99 & 0.28 & Excellent \\
06 Treasury & 100 & $>$0.99 & 0.28 & Excellent \\
07 Inter-DAO & 100 & $>$0.99 & 0.28 & Excellent \\
08 Scale & 100 & $>$0.99 & 0.28 & Excellent \\
09 Voting & 100 & $>$0.99 & 0.28 & Excellent \\
\bottomrule
\end{tabular}
\end{table}

With 100 runs per configuration, all experiments achieve power $>$0.99 for detecting medium effects ($d = 0.5$) and can reliably detect small effects ($d \geq 0.28$).

\subsection{Descriptive Statistics}

\subsubsection{Quorum Sensitivity Experiment}

Table~\ref{tab:quorum-descriptive} reports descriptive statistics for the quorum sensitivity experiment (Experiment 03).

\begin{table}[htbp]
\centering
\caption{Descriptive statistics for quorum sensitivity experiment (Pass Rate)}
\label{tab:quorum-descriptive}
\begin{tabular}{lrrrrr}
\toprule
Quorum (\%) & N & Mean & SD & Min & Max \\
\midrule
1 & 100 & 0.875 & 0.026 & 0.81 & 0.94 \\
3 & 100 & 0.878 & 0.025 & 0.82 & 0.94 \\
5 & 100 & 0.875 & 0.026 & 0.81 & 0.93 \\
8 & 100 & 0.876 & 0.025 & 0.82 & 0.94 \\
10 & 100 & 0.875 & 0.027 & 0.80 & 0.94 \\
\bottomrule
\end{tabular}
\end{table}

\subsection{Regression Results}

\subsubsection{Quorum vs Pass Rate}

Linear regression of proposal pass rate on quorum threshold yields negligible explanatory power:

\begin{table}[htbp]
\centering
\caption{Linear Regression: Pass Rate $\sim$ Quorum}
\label{tab:quorum-regression}
\begin{tabular}{lr}
\toprule
Statistic & Value \\
\midrule
Coefficient & $-0.001$ (SE: 0.003) \\
Intercept & 0.877 (SE: 0.002) \\
$R^2$ & 0.002 \\
Adjusted $R^2$ & $-0.000$ \\
$F$-statistic & 0.89 ($p = 0.35$) \\
$N$ observations & 500 \\
\bottomrule
\end{tabular}
\end{table}

The near-zero $R^2$ confirms that quorum thresholds within the 1--10\% range have negligible effect on pass rates.

\subsection{ANOVA Tables}

\subsubsection{Voting Mechanism Comparison}

One-way ANOVA comparing pass rates across majority, quadratic, and conviction voting:

\begin{table}[htbp]
\centering
\caption{ANOVA: Voting mechanism comparison (Pass Rate)}
\label{tab:anova-voting}
\begin{tabular}{lrrrrr}
\toprule
Source & SS & df & MS & F & p \\
\midrule
Between & 0.002 & 2 & 0.001 & 0.31 & 0.73 \\
Within & 0.391 & 597 & 0.001 & & \\
Total & 0.393 & 599 & & & \\
\bottomrule
\end{tabular}
\end{table}

Effect sizes: $\eta^2 = 0.001$, $\omega^2 < 0.001$ (negligible).

\subsubsection{Scale Effect on Pass Rate}

One-way ANOVA comparing pass rates across DAO sizes (50, 100, 200, 350, 500 members):

\begin{table}[htbp]
\centering
\caption{ANOVA: Scale effect on pass rate}
\label{tab:anova-scale}
\begin{tabular}{lrrrrr}
\toprule
Source & SS & df & MS & F & p \\
\midrule
Between & 8.42 & 4 & 2.105 & 312.4 & $<$0.001 \\
Within & 3.34 & 495 & 0.007 & & \\
Total & 11.76 & 499 & & & \\
\bottomrule
\end{tabular}
\end{table}

Effect sizes: $\eta^2 = 0.72$, $\omega^2 = 0.71$ (very large). Scale has a substantial effect on governance outcomes.

\subsection{Pairwise Comparisons}

\subsubsection{Quorum Sensitivity}

Following non-significant omnibus ANOVA, pairwise comparisons confirm no meaningful differences:

\begin{table}[htbp]
\centering
\caption{Pairwise comparisons for quorum levels (BH-corrected)}
\label{tab:quorum-pairwise}
\begin{tabular}{lcccc}
\toprule
Comparison & Mean Diff & Raw $p$ & Adj $p$ & Significant \\
\midrule
1\% vs 3\% & $-0.003$ & 0.370 & 1.000 & No \\
1\% vs 5\% & $0.000$ & 0.848 & 1.000 & No \\
1\% vs 8\% & $-0.001$ & 0.700 & 1.000 & No \\
1\% vs 10\% & $0.000$ & 1.000 & 1.000 & No \\
3\% vs 5\% & $0.003$ & 0.412 & 1.000 & No \\
\bottomrule
\end{tabular}
\end{table}

All effect sizes are negligible ($|d| < 0.15$).

\subsubsection{Scale Effect Pairwise}

Following significant omnibus ANOVA, pairwise comparisons identify scale thresholds:

\begin{table}[htbp]
\centering
\caption{Pairwise comparisons for DAO scale (BH-corrected)}
\label{tab:scale-pairwise}
\begin{tabular}{lcccc}
\toprule
Comparison & Mean Diff & Cohen's $d$ & Adj $p$ & Significant \\
\midrule
50 vs 100 & $-0.120$ & 1.38 & $<$0.001 & Yes \\
50 vs 200 & $-0.249$ & 2.86 & $<$0.001 & Yes \\
100 vs 200 & $-0.129$ & 1.48 & $<$0.001 & Yes \\
200 vs 350 & $-0.073$ & 0.84 & $<$0.001 & Yes \\
350 vs 500 & $-0.025$ & 0.29 & 0.012 & Yes \\
\bottomrule
\end{tabular}
\end{table}

All scale transitions show statistically significant and practically meaningful differences.

\subsection{Effect Sizes}

\subsubsection{Governance Capture Mitigations}

Effect sizes for vote power cap intervention:

\begin{table}[htbp]
\centering
\caption{Effect sizes for governance capture mitigations}
\label{tab:capture-effects}
\begin{tabular}{lcc}
\toprule
Metric & Cohen's $d$ & Interpretation \\
\midrule
Whale Influence (no cap vs 10\% cap) & 1.42 & Large \\
Capture Risk (no cap vs 10\% cap) & 1.38 & Large \\
Pass Rate (no cap vs 10\% cap) & 0.82 & Large \\
Collusion Threshold (no cap vs 10\% cap) & 0.65 & Medium \\
\bottomrule
\end{tabular}
\end{table}

Vote power caps produce large, practically significant improvements across all governance health metrics.

\subsubsection{Inter-DAO Cooperation}

\begin{table}[htbp]
\centering
\caption{Effect sizes for inter-DAO vs intra-DAO governance}
\label{tab:interdao-effects}
\begin{tabular}{lcc}
\toprule
Comparison & Cohen's $d$ & Interpretation \\
\midrule
Pass Rate (inter vs intra) & 2.31 & Very Large \\
Voting Participation & 0.42 & Small-Medium \\
Time to Decision & 0.89 & Large \\
\bottomrule
\end{tabular}
\end{table}

The large effect size for pass rate confirms that inter-DAO coordination introduces substantial friction.

\subsection{Robustness Checks}

\subsubsection{Normality Assessment}

Shapiro-Wilk tests for normality (selected metrics, N=100):

\begin{table}[htbp]
\centering
\caption{Normality tests for key metrics}
\label{tab:normality}
\begin{tabular}{lccc}
\toprule
Metric & W statistic & $p$-value & Normal? \\
\midrule
Proposal Pass Rate & 0.987 & 0.42 & Yes \\
Average Turnout & 0.991 & 0.71 & Yes \\
Token Gini & 0.985 & 0.33 & Yes \\
Whale Influence & 0.989 & 0.58 & Yes \\
\bottomrule
\end{tabular}
\end{table}

All primary metrics satisfy normality assumptions, supporting use of parametric tests.

\subsubsection{Variance Homogeneity}

Levene's test for homogeneity of variances across experimental groups:

\begin{table}[htbp]
\centering
\caption{Variance homogeneity tests}
\label{tab:levene}
\begin{tabular}{lccc}
\toprule
Experiment & Levene $F$ & $p$-value & Equal Variance? \\
\midrule
Quorum Sensitivity & 0.89 & 0.47 & Yes \\
Voting Mechanisms & 0.67 & 0.51 & Yes \\
Scale Sweep & 2.34 & 0.054 & Yes (marginal) \\
\bottomrule
\end{tabular}
\end{table}

Welch's $t$-tests were used for pairwise comparisons to account for potential heteroscedasticity.
